\scanpage{87}
\mspage{93}

Mais on peut s'inspirer de cet exemple pour le modifier de façon qu'il se transpose à des ens. ordonnés. Il suffit de réaliser, pour $\unsure{G|C}$ donné

a) Une opération \deleted{\unsure{sans point fixe}}\inserted{\uline{libre} de $G$} sur un \deleted{catégorie}\inserted{objet} $E$ de $\mathit{Cat}$, \deleted{qui soit} \inserted{telle que} \deleted{$\mathbb D$-asphérique (i.e. à la cohomologie triviale)} $H^{\unsure{\cdot|*}}(E,k)^G\simeq k$, i.e. $B=E/G$ soit $\mathbb D$-asphérique

b) Une opération de $G$ sur un objet $Z$, telle que l'inclusion

\[
H^{\cdot}(Z,k)^G\unsure{\relbar\joinrel\relbar|\longrightarrow}H^{\cdot}(Z,k)^G
\]

ne soit pas iso.

c) \deleted{\unsure{D'}}\inserted{Choisir pour $E,Z$ des ens. ordonnés, et} S'assurer que $B=E/G$ est également représenté par une catégorie ordonnée. Sauf erreur, il en sera de $=$ pour $X=E\times Z/G$.

\unsure{Donc|D'où} on aura en tous les cas

\diagram{\begin{tikzcd}[ampersand replacement=\&,row sep=large,column sep=large]
X\arrow[d] \& \overset{\substack{Z\times E\\ \shortparallel}}{X'}\arrow[l]\arrow[d] \& Z\arrow[l]\arrow[d,no head]\\
\underset{\substack{=E/G\\ \text{asph.}}}{S}\&\underset{\substack{=E\\ \text{asph.}\\ \text{si on peut}}}{S'}\arrow[l]\&
\underset{\text{asph.}}{S''=\{\underset{\substack{\rotatebox{-90}{$\in$}\\ E}}s\}}\arrow[l]
\end{tikzcd}}

Cela donne encore un exemple avec $X$ $\mathbb D$-fibré propre et lisse, sur une base $S$ $\mathbb D$-asphérique, telle que l'inclusion d'une fibre $X_s$ dans $X$ ne soit pas une $\mathbb D$-équivalence. On peut $=$ se débrouiller pour que $H^{\cdot}(Z,k)^G\simeq k$, ce qui \unsure{implique|signifie} que

% archive sheet number: 86
